本节将传播参数的拟合与感染阈值的设定分开。西安市 2021 年末常住人口为 $N=13{,}163{,}000$\cite{XianStat2021}。现有资料尚不足以确定与本模型全部新增感染口径一致的 ICU 需求函数、初始病例未来负担及同期可用容量，因而本节没有按第\ref{sec:model:strategy}节的充分条件校准阈值。以下保留一个文献比例参考及其附近的设计情景，用于比较传播与干预指标。

为与已有情景保持一致，仅在本小节使用参考系数 $\rho_{\mathrm{ICU}}>0$，按形式关系 $\eta=C/\rho_{\mathrm{ICU}}$、$\theta=(C/N)/\rho_{\mathrm{ICU}}$ 说明阈值量级。陈胤孜等\cite{Chen2021ICUForecast}基于既往资源资料预测，2021 年全国每10万人综合 ICU 床位约为 $4.37$ 张；它不是西安同期可用容量的实测值。Angulo 等\cite[补充材料 S4.1]{Angulo2021}对早期 COVID-19 采用需要重症监护的感染者比例名义值 $0.60\times0.15\times0.28=0.0252$。暂将全国预测量作为人均容量参考、不扣除既有占用，并把该名义比例移用于本节的参考系数，得到
\[
\theta=\frac{4.37/10^5}{0.0252}\approx1.73\times10^{-3},
\qquad \eta\approx2.28\times10^4.
\]
这一计算只给出比例情景的量级。$\rho_{\mathrm{ICU}}$ 没有依据共同需求函数及初始负担估计，不是第\ref{sec:model:strategy}节充分条件中的 $pL\beta c_0S_0/N$，也不是所有策略下实际占用与 $I(t)$ 的固定比值。以确诊病例或全部感染者为分母得到的 ICU 入院风险本身就可能不同\cite{Pung2023Undetected}，上述移植不能作为本地临床校准。

为比较这一量级附近的控制结果，保留代表性设计情景
\begin{equation}
\eta=0.002N=26{,}326.
\end{equation}
这不是上述名义换算值的保守取整。在 $\rho_{\mathrm{ICU}}=0.0252$ 的比例情景中，它对应每10万人 $5.04$ 张、全市约 $663$ 张的名义容量；这些数值不作为西安实际可用资源。第\ref{sec:joint:compare-numerics}节及本节采用该阈值的算例，均为同一社区感染控制目标下的策略比较，尚未核查第\ref{sec:model:strategy}节的容量充分条件。它们既不构成实际 ICU 不超容的证明，也不能据未校准而判定必然超容。阈值变化的影响仍由后续敏感性分析说明。
